\scanpage{34}
\uline{Catégorie des algèbres bivariantes} (sur un anneau $k$)

C'est la catégorie dont les objets sont les $k$-algèbres \inserted{\uline{muni d'une forme} $\unsure{k|\kappa}_A:A\mathbin{\text{---}}k$}, et les morphismes de $A$ dans $B$ les couples $f=(f_*,f^*)$, où $f^*:B\mathbin{\text{---}}A$ est un hom d'algèbres, et où $f_*:A\mathbin{\text{---}}B$ est un hom $B$ linéaire \inserted{(à g. et à dr.)}, i.e. un hom additif satisfaisant les ``form. de projection''
\[
(*)\qquad[f_*(xf^*(y))=f_*(x)y],\quad f_*(f^*(y)x)=yf_*(x).
\]
\margin{(signe $(-1)^{\delta j}$ \unsure{raisonnable} ds le cas gradué; si la deux. formule est valable, $y$ de degré $j$, $f_*$ de degré $\delta$.)}
\margin{Il faut choisir entre l'une et l'autre formule!}
et de plus:
\[
(**)\qquad\unsure{k|\kappa}_Bf_*=\unsure{k|\kappa}_A.
\]
On a une variante gr.duée, \deleted{et} On aura à travailler dans diverses sous-catégories pleines, ou même non pleines.

Dans le cas gradué anticommutatif (où $f^*$ est donc un hom d'Alg. gr.duées, et \unsure{où} $f_*$ préserve à translation \inserted{\uline{de degré pair}} près la graduation), les deux formules $(*)$ sont équivalentes.

\margin{\uline{Remarque} \unsure{Bien que|Pour que} $(f_*,f^*)$ soit un isomorphisme, \unsure{il faut que $f^*$ et} \deleted{\unsure{un isom de}} $f_*$ \unsure{bijectifs}, i.e. $f^*$ un isom et $f_*(1)$ inversible; les ``bons'' isom. sont ceux où $f_*(1)=1$, ou encore $f_*$ et $f^*$ inverses l'un de l'autre.}
\figure{\begin{tikzpicture}\draw(0,5)..controls(1,6)and(1.7,5.8)..(1.8,4.5)..controls(1.8,4)and(2.3,4.3)..(2.7,4)..controls(3,3.7)and(2.5,3.6)..(2.5,3.1)--(2.4,1.2)..controls(2.2,.6)and(1.5,.3)..(.3,0);\end{tikzpicture}}

Soit $A$ une algèbre muni d'une forme
\[
\unsure{k|\kappa}_A:A\mathbin{\text{---}}k
\]
\deleted{que}\inserted{donc} la forme bilinéaire
\[
(x,y)\mapsto\unsure{k|\kappa}(xy)
\]
sur $A\times A$ \deleted{est séparée, donc} définit un hom. \inserted{\unsure{(= dual!)}} \deleted{injectif}
\[
\varphi_A:A\longrightarrow\check A\qquad(\check A:\mathrm{Hom}_{k\text{-}\mathrm{mod}}(A,k)).
\]
\inserted{$\langle y,\varphi_A(x)\rangle=\unsure{k|\kappa}(yx)$}
\margin{\deleted{\unsure{l'isomorphisme $\mathcal V(B)\mathbin{\text{---}}\mathcal V(A)$}}}
Supposons que $B$ est de même muni de $\unsure{k|\kappa}_B:B\mathbin{\text{---}}k$, \inserted{et que $\varphi_B$ soit \uline{injectif}}, et soit $f^*:B\longrightarrow A$ tel que l'hom transposé ${}^tf^*:\check A\mathbin{\text{---}}\check B$ \unsure{applique|vérifie} $\mathrm{Im}\,\varphi_A$ dans $\mathrm{Im}\,\varphi_B$ [ce qui n'est pas une restriction si $\varphi_B:B\overset{\sim}\longrightarrow\check B$!].

Alors l'application induite, \inserted{notée} $f_*$, est donc caractérisée par
\[
(***)\qquad\unsure{k|\kappa}(yf_*(x))=\unsure{k|\kappa}(f^*(y)x)
\]
et satisfait de ce fait la $2^{\text{ième}}$ relation $(*)$, [donc la première dans le cas commutatif, ou gradué anticommutatif, quand $f_*$ est hom à translation \inserted{paire} près].
\inserted{De plus, il satisfait $(**)$ (faire $y=1$ dans $(***)$).}
\deleted{\unsure{Supposons donc $A,B$}}

% Addition on the lowest edge of the sheet.
\inserted{Réciproquement \unsure{étant donné un hom} $(f^*,f_*)$ de $(A,\unsure{k|\kappa}_A)$ dans $(B,\unsure{k|\kappa}_B)$, $f_*$ satisfait $(***)$ ce qui \unsure{le} détermine par $f^*$. \deleted{\unsure{$\unsure{k|\kappa}_A$ dans $\unsure{k|\kappa}_B$, de $(**)$ donc $f_*$ est}}}
\mspage{33}
